\section{Risk-neutral applications beyond the original terminal basket}
\label{sec:moment-applications}

The preceding moment theorem concerns a mathematical class, rather than an
extension obtained by substituting new data into the original GBM. Transform
controls and conditional Monte Carlo across general financial models already
have published antecedents \citep{dingec}. I distinguish those established
methods from the finite-moment and encoding consequences investigated here.

\begin{theorem}[A certificate without an affine or Gaussian return model]
\label{thm:positive-asset-moment}
Let $S_1,\ldots,S_n$ be strictly positive random terminal asset values
under a stipulated pricing law, with $\E S_i^2<\infty$. Let positive
weights sum to one and set $A=\sum_iw_iS_i$, $G=\prod_iS_i^{w_i}$.
Let $F:\mathbb R_+\to[0,H]$ be increasing and $c_F$-Lipschitz,
and let $\mathcal H$ be any sub-$\sigma$-algebra. Define
\begin{equation}
 X_0=\{F(A)-F(G)\}/H,\qquad
 X=\E[X_0\mid\mathcal H],\qquad
 \Delta=\E[A-G\mid\mathcal H].
\end{equation}
Then $X\in[0,1]$ and
\begin{equation}
 \E X^2\leq\min\{1,c_F^2\E\Delta^2/H^2\}
 \leq\min\{1,c_F^2\E(A-G)^2/H^2\}.
\end{equation}
For a finite-law implementation $\widehat X\in[0,1]$ with a
proved pointwise error at most $\eta$, its variance is bounded by
\begin{equation}
 \operatorname{Var}(\widehat X)
 \leq\min\{1/4,(\sqrt{\min\{1,c_F^2\E\Delta^2/H^2\}}+\eta)^2\},
\end{equation}
provided the expectation in this last expression concerns the same
ideal finite law. A continuous-law expectation cannot be substituted
without a separate distribution-transfer argument.
\end{theorem}

\begin{proof}
Weighted arithmetic--geometric ordering gives $A\geq G>0$.
The payoff conditions yield
$0\leq X_0\leq\min\{1,c_F(A-G)/H\}$.
Conditional expectation preserves these inequalities, so
$0\leq X\leq\min\{1,c_F\Delta/H\}$.
The square-moment assertions follow by squaring and then by conditional
Jensen's inequality, $\Delta^2\leq\E[(A-G)^2\mid\mathcal H]$.
All expectations exist: convexity gives
$A^2\leq\sum_iw_iS_i^2$, while $0\leq A-G\leq A$.
The pointwise implementation assumption gives
$\|\widehat X\|_2\leq\|X\|_2+\eta$ by the triangle inequality
in $L^2$. Subtracting the squared mean and using the bounded-range
variance bound $1/4$ proves the last assertion.
\end{proof}

For the capped call spread $c_F=1$. The unconditioned case takes
$\mathcal H$ to contain all terminal asset information; a tractable
conditioning factor can reduce the square moment further. This statement
does not require a log-affine model, independence or a moment-generating
function. Its practical use still requires a computable conservative
moment certificate and a checked numerical law. The subsequent transform
formula is one such route, rather than an assertion that all financial
models satisfy its moment assumptions.

\begin{proposition}[The correct general risk-neutral shift]
\label{prop:mgf-risk-neutral}
Let $\psi(v)=\log\Lambda(v)$ be finite and twice continuously differentiable
on a neighbourhood of every $R_0+\varepsilon b_i$ for
$|\varepsilon|\leq\varepsilon_0$, with its Hessian bounded on those
compact line segments. Let $Z\sim N(0,1)$ be independent of $Y$ and define
the rebased terminal assets
\begin{equation}
 S_{i,\varepsilon}=K\exp\{\mu_i(\varepsilon)
                  +(R_0+\varepsilon b_i)\cdot Y+\lambda Z\},\qquad
 \mu_i(\varepsilon)=r-d-\lambda^2/2-\psi(R_0+\varepsilon b_i).
\end{equation}
Then $\E S_{i,\varepsilon}=Ke^{r-d}$.
Suppose $w_i>0$, $\sum_iw_i=1$, $\sum_iw_i b_i=0$, and the
integrability hypothesis of
Theorem~\ref{thm:general-weak-dispersion} holds for
$V=Y-\nabla\psi(R_0)$. Let $P:\R\to[0,1]$ be increasing,
continuously differentiable and have a globally bounded derivative. The
leading fractional residual is
\begin{equation}
 \frac{X_\varepsilon}{\varepsilon^2}
 \longrightarrow
 \tfrac12P'(m_0(Y))\sum_iw_i
       \{b_i\cdot(Y-\nabla\psi(R_0))\}^2
 \quad\hbox{in }L^2,
\end{equation}
where $m_0(Y)=r-d-\lambda^2/2-\psi(R_0)+R_0\cdot Y$.
\end{proposition}

\begin{proof}
Independence and the normal exponential moment give
$\E S_{i,\varepsilon}
=Ke^{\mu_i(\varepsilon)}\Lambda(R_0+\varepsilon b_i)
e^{\lambda^2/2}=Ke^{r-d}$.
Taylor's integral formula writes
\begin{equation}
 \psi(R_0+\varepsilon b_i)
 =\psi(R_0)+\varepsilon\nabla\psi(R_0)\cdot b_i
       +\varepsilon^2 e_i(\varepsilon),\qquad
 e_i(\varepsilon)=\int_0^1(1-t)
 b_i^{\mathsf T}\nabla^2\psi(R_0+t\varepsilon b_i)b_i\,dt.
\end{equation}
The Hessian hypothesis bounds every $e_i$ uniformly. Subtracting the
weighted geometric log gives the exact relative logs
\begin{equation}
 D_{i,\varepsilon}
 =\varepsilon b_i\cdot(Y-\nabla\psi(R_0))
 -\varepsilon^2\{e_i(\varepsilon)-\sum_kw_ke_k(\varepsilon)\}.
\end{equation}
Their intercepts have weighted mean zero and are uniformly
$O(\varepsilon^2)$. The geometric base converges almost surely to
$m_0(Y)$ because the weighted linear Taylor term is zero.
The hypotheses therefore match
Theorem~\ref{thm:general-weak-dispersion}; its $L^2$ conclusion proves
the assertion. No Gaussian drift is reused after changing the latent law.
\end{proof}

This drift calculation establishes consistency of the stipulated pricing
measure. It does not calibrate that measure or validate an observed exotic
price. Rebased components allow the moment theorem to concern portfolio
returns, while its unrestricted affine constants also accommodate unequal
initial asset levels in the non-asymptotic result. The weak-dispersion
specialisation requires the stated common limiting base.

\subsection{A variance-gamma financial example and its moment domain}

The symmetric variance-gamma construction is established financial
modelling \citep{madan}. I use $Y=\sqrt G\,N$, where $N\sim N(0,1)$,
$G$ is independent with gamma shape $1/\nu$ and scale $\nu$, and
$\nu=0.4$. I derive the required terminal moments, rather than transfer
an empirical claim from that source. With $k=1/\nu$, the gamma density is
$g^{k-1}e^{-g/\nu}/\{\Gamma(k)\nu^k\}$ for $g>0$. Conditional on $G$,
$\E(e^{tY}\mid G)=e^{t^2G/2}$. Its integral is
\begin{equation}
 M_Y(t)=\frac1{\Gamma(k)\nu^k}
 \int_0^\infty g^{k-1}e^{-(1/\nu-t^2/2)g}\,dg
 =(1-\nu t^2/2)^{-1/\nu},\qquad |t|<\sqrt{2/\nu}.
\end{equation}
The equality follows by changing variables in the defining gamma
integral; outside the stated open domain the necessary integral is
not finite. Gamma recursion gives $\E G=1$ and
$\E G^2=1+\nu$. Conditioning on $G$ therefore gives
$\E Y^2=1$ and $\E Y^4=3(1+\nu)=4.2$.

For equal weights and $b_1=-0.2,b_2=0.2$, $R_0=0$,
$\varepsilon_0=1$, the norm-exponential constant is
$2\varepsilon_0B=0.4$. For any $\delta>0$ with
$0.4+\delta<\sqrt5$,
$\E e^{(0.4+\delta)|Y|}\leq
M_Y(0.4+\delta)+M_Y(-0.4-\delta)<\infty$.
The exponential series gives
$|Y|^4\leq4!\delta^{-4}e^{\delta|Y|}$, so the exact
polynomial-times-exponential hypothesis is satisfied.
The drift is $r-d-\lambda^2/2-\log M_Y(0.2\varepsilon)$.
This is a correctly normalised non-Gaussian terminal pricing family;
the original fixed Gaussian drifts would give a different model.

In this symmetric family the relative log ratio is
$J_\varepsilon=\log\cosh(0.2\varepsilon Y)$ and the geometric base
is deterministic before integrating the independent common normal shock.
Writing $a=P'(r-d-\lambda^2/2)>0$, the exact limit coefficients are
\begin{equation}
 \E H_0=0.02a,\qquad
 \operatorname{Var}(H_0)=0.0004a^2\operatorname{Var}(Y^2).
\end{equation}
The variance factors are $2$ for a standard normal, $0.8$ for a
standardised uniform law, $3.2$ for this gamma-normal mixture, and zero
for a symmetric two-point law. In the last case $Y^2=1$ and the
entire conditional fractional payoff is deterministic, not just its
leading coefficient. A flag drawn at that constant probability can still
have positive variance. This is a boundary of the estimator comparison;
it is not a lower bound on a quantum method that uses the full value oracle.

\subsection{Monitored Asian baskets and a common Gaussian factor}

The general affine-log-price theorem can also cover a finite collection
of monitored GBM prices with positive averaging weights. Geometric
average controls and quantum Asian valuation are established
\citep{kemna,rebentrost}. Let their non-singular joint Gaussian
log-covariance be $\mathcal C$. For a common load $\lambda$,
\begin{equation}\label{eq:common-factor-domain}
 \mathcal C-\lambda^2\mathbf1\mathbf1^{\mathsf T}\succeq0
 \quad\Longleftrightarrow\quad
 \lambda^2\leq(\mathbf1^{\mathsf T}\mathcal C^{-1}\mathbf1)^{-1}.
\end{equation}
Indeed, Cauchy--Schwarz gives
$(\mathbf1^{\mathsf T}x)^2\leq
(\mathbf1^{\mathsf T}\mathcal C^{-1}\mathbf1)(x^{\mathsf T}\mathcal Cx)$,
proving sufficiency. Necessity follows by substituting
$x=\mathcal C^{-1}\mathbf1$ into the quadratic-form inequality.
A square root of the residual covariance then constructs independent
$Y,Z$ with log prices $a_i+B_i\cdot Y+\lambda Z$.
Thus the moment theorem applies to the resulting arithmetic/geometric
monitored average once its kernel and moments are supplied.
Singular monitoring covariances require a separate range condition;
I do not assert~\eqref{eq:common-factor-domain} using an inverse where
none exists. This application states a mathematical scope beyond the
original terminal experiments, rather than claiming to have executed
new Asian circuits or demonstrated scalability on market data.
